\chapter{Step 08: A Program Loader}
\label{ch:step08}

A computer needs memory, and the 512 bytes of a block RAM are not much. The
iCE40UP5K has four larger memory blocks, together 128\,KiB, and in this step we
put them to use: as the 64-bit memory of the ULM, which can be read and written
in bytes, words (2 bytes), longs (4 bytes) and quads (8 bytes). To get
something into the memory, we build a \term{loader}: it reads hexadecimal
digits from the I/O device of step 07, combines two of them into a byte, and
stores the bytes one after the other, until the character \cc{0x04} ends the
input. A small debugger lets us look at the result with the buttons. New are
the SPRAM blocks of the FPGA and how we instantiate them, a simulation model
that the synthesis must not see, the shift and mask logic for accesses of
different sizes, and \cc{for} loops. Everything is in \cc{08_dev_loader}; the
I/O device and the FIFOs come from step 07 (Chapter~\ref{ch:step07}).

\section{What the Design Does}
\label{sec:step08-design}

In a terminal connected to the board you type a program as hexadecimal
digits, for example \cc{32 10}, and then Ctrl-D, which sends the character
\cc{0x04}. The loader stores \cc{0x32} at address 0 and \cc{0x10} at address 1.
Characters that are not hexadecimal digits, like the space, are skipped. While
you type, the display shows the byte being assembled, and \cc{LED4} is on
while half of a byte (one digit) has been typed. After \cc{0x04}, \cc{LED3}
goes on and the memory belongs to the debugger: \cc{BTN1} reads the byte at
address 0, each press of \cc{BTN2} the byte at the next address, and the
display shows the low byte of the address, or the byte read while \cc{BTN3}
is held.
Pressing \cc{BTN1} and \cc{BTN3} together resets the loader.

Figure~\ref{fig:step08-blocks} shows how the parts are connected. Three
interfaces of type \cc{if_ram} connect the loader and the debugger to a switch,
and the switch to the memory \cc{dev_ram}. The switch is controlled by the
loader's output \cc{done}.

\begin{figure}[ht]
\centering
\begin{tikzpicture}[node distance=5mm and 8mm]
  \node[box,minimum width=16mm,minimum height=9mm] (io) {\cc{dev_io}};
  \node[box,right=of io,minimum width=30mm,minimum height=9mm] (ld) {\cc{dev_loader}};
  \node[box,below=10mm of ld,minimum width=30mm,minimum height=9mm] (db) {\cc{dev_ram_debugger}};
  \node[box,right=24mm of ld,minimum width=18mm,minimum height=22mm,yshift=-9.5mm] (sw) {\cc{dev_ram_switch}};
  \node[box,fill=green!10,right=of sw,minimum width=18mm,minimum height=22mm] (ram) {\cc{dev_ram}\\[-2pt]\footnotesize 4 SPRAMs\\[-2pt]\footnotesize 128\,KiB};
  \node[left=8mm of io,note] (rx) {\cc{RX}, \cc{TX}};
  \draw[biarrow] (rx) -- (io);
  \draw[biarrow] (io) -- node[above,note] {\cc{if_io}} (ld);
  \draw[biarrow] (ld.east) -- node[above,note] {\cc{ram_loader}} (ld.east -| sw.west);
  \draw[biarrow] (db.east) -- node[above,note] {\cc{ram_debugger}} (db.east -| sw.west);
  \draw[biarrow] (sw) -- node[above,note] {\cc{ram}} (ram);
  \draw[arrow] (ld.north) -- ++(0,5mm) -| node[pos=0.25,above,note] {\cc{done} = select} (sw.north);
  \node[left=8mm of db,note] (btn) {\cc{BTN1}--\cc{BTN3}};
  \draw[arrow] (btn) -- (db);
\end{tikzpicture}
\caption{The parts of step 08. Until the loader is done, the switch connects
it to the memory; afterwards, the debugger.}
\label{fig:step08-blocks}
\end{figure}

\section{The Single-Port RAMs of the iCE40UP5K}
\label{sec:step08-spram}

Besides the 30 block RAMs of Section~\ref{sec:step07-bram}, the iCE40UP5K has
four \term{single-port RAMs}, SPRAM for short. Each has $256$\,kbit, organised
as $16384$ words of 16 bits, so together they hold $4 \cdot 256\,\text{kbit} =
1\,\text{Mbit} = 128$\,KiB, more than eight times as much as all block RAMs
($30 \cdot 4\,\text{kbit} = 15$\,KiB). The price is in the name: a single port
means one address per clock cycle, used either for reading or for writing.

Instead of describing a memory and hoping that yosys recognises it, as in
step 07, we instantiate the hard block directly, like a module. Its name in the netlist is
\cc{SB_SPRAM256KA}, and \cc{spram.sv} wraps it:

\svfile{08_dev_loader/spram.sv}

There is no module \cc{SB_SPRAM256KA} in our design that describes how it
works. yosys knows it as one of the building blocks of the iCE40, like
\cc{SB_LUT4}, and puts it into the netlist as it is; nextpnr places it on one of
the four SPRAM sites. Its ports:
\begin{description}
\item[\cc{ADDRESS}, \cc{DATAIN}, \cc{DATAOUT}:] 14 address bits ($2^{14} =
  16384$ words) and 16 data bits in each direction.
\item[\cc{WREN}, \cc{MASKWREN}:] \cc{WREN} = 1 writes, \cc{WREN} = 0 reads.
  \cc{MASKWREN} has one bit per 4-bit group (\term{nibble}) of the word, and
  only the nibbles whose bit is 1 are written. So we can write a single byte of
  a word by setting two of the four bits. \cc{spram.sv} uses one signal
  \cc{we} for both: write if any mask bit is set.
\item[\cc{CLOCK}:] both reading and writing happen at the rising edge. A read
  delivers the word at \cc{DATAOUT} after the edge, like the block RAM of step
  07.
\item[\cc{CHIPSELECT}, \cc{STANDBY}, \cc{SLEEP}, \cc{POWEROFF}:] for power
  saving. We keep the block always selected and awake. Note that
  \cc{POWEROFF} is active low: 1 means ``power on''.
\end{description}

\subsection*{A model for the simulation, hidden from the synthesis}

Verilator does not know \cc{SB_SPRAM256KA}: it is not a SystemVerilog module
but a piece of silicon. For the simulation we need a description of what it
does, a \term{simulation model}. The file \cc{SB_SPRAM256KA.sv} contains one,
taken from yosys's own models of the iCE40 cells:

\svpart{firstline=11,lastline=53}{08_dev_loader/SB_SPRAM256KA.sv}

The model is a memory \cc{mem} of 16384 words, here declared with an explicit
range \cc{[0:16383]} instead of the size \cc{[DEPTH]} of step 07; both mean
the same. At the rising clock edge it either reads into \cc{DATAOUT} or writes
the selected nibbles, and while writing \cc{DATAOUT} becomes undefined
(\cc{16'bx}). This is how we learn the behaviour of the hard block without the
board.

The body of the model is enclosed in \cc{`ifndef SYNTHESIS} \dots\
\cc{`endif}. As with the include guard of Section~\ref{sec:step07-package},
the preprocessor drops the text if the macro \cc{SYNTHESIS} is defined. The
Makefile calls sv2v with \cc{--define=SYNTHESIS} (Section~\ref{sec:step01-synth}),
Verilator without. So Verilator simulates the model, while the file that
yosys reads contains only the ports:

\outputfile{step08_sv2v_spram.txt}

An empty module like this is a \term{black box}: it says which connections the
block has, but nothing about its inside, which is exactly right for a hard
block. Without the guard, a tool that does not know \cc{SB_SPRAM256KA} would
have to build a memory of 256\,kbit from logic cells, and the whole chip does
not have that many flip-flops. (yosys's \cc{synth_ice40} happens to be safe
anyway: it loads its own definition of the cell, and we checked that it
produces the same statistics without \cc{--define=SYNTHESIS}. The guard does not
rely on that.)

\section{for Loops Are Replicated Hardware}
\label{sec:step08-for}

The model contains our first \cc{for} loop:
\begin{lstlisting}
always @(negedge POWEROFF) begin
    for (i = 0; i <= 16383; i = i+1)
        mem[i] = 16'bx;
end
\end{lstlisting}
In the simulation this is what it looks like: a loop that the simulator runs,
setting all words to ``undefined'' when the power goes off. It is part of the
model, which the synthesis never sees.

In a description of hardware, a \cc{for} loop means something else. There is no
time in which a loop could run; the tools \term{unroll} it, that is, they write
the body down once for each value of the loop variable, and each copy becomes
its own piece of circuit. The number of iterations must be a constant. An
example: the module \cc{quad_lshift} below (Section~\ref{sec:step08-ram})
selects one of eight shifts with a \cc{case}. The same can be written with a
loop:
\begin{lstlisting}
always_comb begin
    data_out = data_in;
    for (int k = 1; k < 8; k++)
        if (offset == k)
            data_out = data_in << 8*k;
end
\end{lstlisting}
Unrolled, this is seven \cc{if} statements one after the other, and each
\cc{if} is a comparator and a 64-bit multiplexer that either passes on the
value from the copy before or replaces it: a chain of seven multiplexers. In
\cc{always_comb}, a later assignment to the same signal wins, so the chain is
ordered, with the last copy closest to the output. We synthesised the
repository's version, this loop, and a third version with a single shift,
\cc{data_out = data_in << {offset, 3'b000};} (the comment in
\cc{quad_lshift.sv}):

\outputfile{step08_lshift_variants.txt}

All three compute the same function. The loop is a chain of seven 2-to-1
multiplexers and costs the most; the \cc{case} is one 8-to-1 multiplexer per
bit; the single shift by $8 \cdot \mathtt{offset}$ lets yosys build a
\term{barrel shifter}: three stages that shift by 8, 16 and 32 bits or not,
one stage per bit of \cc{offset}. yosys simplifies a lot, but not everything;
how we describe a function still shows in the result.

\section{A 64-bit Memory from Four 16-bit Ones}
\label{sec:step08-ram}

The ULM has 64-bit registers and reads and writes bytes, words, longs and
quads. We place the four SPRAMs side by side, so that each address of the
SPRAMs holds a quad of 64 bits (Figure~\ref{fig:step08-layout}). The memory is
addressed in bytes: a byte address $a$ has 17 bits, since $2^{17} = 128$\,Ki,
and we split it as
\[
  a = 8t + o, \qquad t = \lfloor a/8 \rfloor \in \{0,\dots,16383\}, \quad
  o = a \bmod 8 \in \{0,\dots,7\},
\]
where the \term{tag} $t$ (the upper 14 bits of $a$) is the address of the quad
in all four SPRAMs and the \term{offset} $o$ (the lower 3 bits) selects the
byte within the quad. The byte with offset 0 is the most significant byte of
the quad: the bytes are stored in big-endian order, as on the ULM.

\begin{figure}[ht]
\centering
\begin{tikzpicture}[x=1.6mm,y=1mm]
  \foreach \s/\n in {0/spram00,1/spram01,2/spram10,3/spram11} {
    \draw[fill=green!8] (16*\s,0) rectangle (16*\s+16,8);
    \node[font=\scriptsize\ttfamily] at (16*\s+8,15) {\n};
    \pgfmathtruncatemacro{\hi}{63-16*\s}
    \pgfmathtruncatemacro{\lo}{48-16*\s}
    \node[font=\scriptsize,anchor=south west] at (16*\s,8) {\hi};
    \node[font=\scriptsize,anchor=south east] at (16*\s+16,8) {\lo};
  }
  \foreach \b/\v in {0/01,1/23,2/45,3/67,4/89,5/AB,6/CD,7/EF} {
    \draw (8*\b,0) -- (8*\b,8);
    \node[font=\small\ttfamily] at (8*\b+4,4) {\v};
    \node[font=\scriptsize] at (8*\b+4,-3) {$o=\b$};
  }
  \node[note,anchor=east] at (-2,4) {quad at tag $t$:};
  \node[note,anchor=north] at (32,-7) {byte addresses $8t+0, \dots, 8t+7$};
\end{tikzpicture}
\caption{One quad of the memory, here with the value \cc{0x0123456789ABCDEF}.
Each SPRAM holds 16 of its 64 bits, all four at the same address $t$. The byte
at address $8t+o$ is bits $63-8o$ down to $56-8o$.}
\label{fig:step08-layout}
\end{figure}

\subsection*{The types of an access}

The package \cc{pkg_ram} from Section~\ref{sec:step07-package} defines the
kinds of access:
\svpart{firstline=12,lastline=31}{08_dev_loader/pkg_ram.sv}
and the interface to the memory uses them as types:

\svfile{08_dev_loader/if_ram.sv}

A user sets \cc{op} to \cc{RAM_FETCH} or \cc{RAM_STORE} for one clock cycle,
together with \cc{data_type}, the byte address \cc{addr} and, for a store,
\cc{data_in}. A value of 8, 16 or 32 bits is always passed in the low bits of
the 64-bit \cc{data_in} and \cc{data_out}.

\subsection*{Reading: shift left, then extend}

\svpart{firstline=16,lastline=51}{08_dev_loader/dev_ram.sv}

The SPRAM delivers its data one clock cycle after the address. By then the
user has usually taken \cc{op} back to \cc{RAM_NOP} and may have changed
\cc{addr}. So \cc{dev_ram} stores the address and data type of every access in
the registers \cc{addr_r} and \cc{data_type_r}, and the \cc{always_comb} block
selects: during an access the new values, otherwise the stored ones. Thus the
offset that the shifters see belongs to the data that the SPRAMs deliver. (In
cycles without an access, \cc{we} is 0, and the SPRAMs keep reading the quad at
the stored address.)

The rest of \cc{dev_ram} instantiates the four SPRAMs with the tag as
address, and between them and the interface the units that move the bytes.
The read side:

\svpart{firstline=53,lastline=72}{08_dev_loader/dev_ram.sv}

To read the value at offset $o$, \cc{quad_lshift} shifts the quad left by $8o$
bits, so that the wanted bytes become the most significant ones:

\svfile{08_dev_loader/quad_lshift.sv}

(\cc{<<} shifts left and fills with zeros; the bits that move out at the top are
lost. The imports of \cc{RAM_LONG} and so on are not used here.) Then
\cc{quad_ext_high} takes the top 8, 16, 32 or 64 bits and puts them into the
low bits of the result, filling the rest with zeros:

\svpart{firstline=21,lastline=41}{08_dev_loader/quad_ext_high.sv}

\cc{{RAM_QUAD_SIZE-RAM_BYTE_SIZE{1'b0}}} is a replication (Chapter~\ref{ch:step05}): 56 copies of
the bit 0, concatenated. So, for a byte, \cc{data_out} is 56 zeros followed by
bits 63 to 56 of the shifted quad. The \cc{default} case, here for
\cc{RAM_QUAD}, passes everything.

With the quad of Figure~\ref{fig:step08-layout}, a byte fetch at offset 3:
\[
\begin{array}{lll}
\text{from the SPRAMs} & \mathtt{0x0123456789ABCDEF} & \\
\text{shifted left by } 8 \cdot 3 = 24 & \mathtt{0x6789ABCDEF000000} &
  \text{(\cc{quad_lshift})} \\
\text{top 8 bits, zero-extended} & \mathtt{0x0000000000000067} &
  \text{(\cc{quad_ext_high})}
\end{array}
\]
and a long fetch at offset 4 gives \cc{0x89ABCDEF00000000} after the shift and
\cc{0x0000000089ABCDEF} at the end.

\subsection*{Writing: extend, shift right, mask}

A store runs the other way. \cc{quad_ext_low} moves the low 8, 16 or 32 bits of
\cc{data_in} to the top of a quad, with zeros below (it is the mirror image of
\cc{quad_ext_high}), and \cc{quad_rshift} shifts right by $8o$, which moves
them to their place. Then all four SPRAMs get 16 bits each, but only the
nibbles that belong to the value may be written. This is the job of
\cc{quad_we_mask}:

\svpart{firstline=23,lastline=54}{08_dev_loader/quad_we_mask.sv}

The 16 mask bits \cc{{we00, we01, we10, we11}} correspond to the 16 nibbles of
the quad, the first bit to bits 63--60. A byte is two nibbles, so a byte store
at offset $o$ sets bits $2o$ and $2o+1$ counted from the left: \cc{16'hC000}
for offset 0, \cc{16'h3000} for offset 1, and so on. Here is a byte store of
\cc{0xAB} at offset 5:
\[
\begin{array}{lll}
\text{\cc{data_in}} & \mathtt{0x00000000000000AB} & \\
\text{to the top} & \mathtt{0xAB00000000000000} & \text{(\cc{quad_ext_low})} \\
\text{shifted right by } 8 \cdot 5 = 40 & \mathtt{0x0000000000AB0000} &
  \text{(\cc{quad_rshift})} \\
\text{mask} & \mathtt{0x0030} = 0000\;0000\;0011\;0000_2 &
  \text{(\cc{quad_we_mask})}
\end{array}
\]
The mask enables two nibbles of \cc{spram10}, and there the 16 bits are
\cc{0x00AB}: exactly the byte \cc{0xAB} lands in bits 23 to 16 of the quad,
and the other seven bytes in memory stay as they were.

For words and longs, the mask is set only if the offset is a multiple of the
size. A store to an unaligned address, say a word at offset 1, writes
nothing.

\begin{exercise}
\label{ex:step08-fetch}
The quad at tag $t$ is \cc{0x0123456789ABCDEF}. What does a word fetch from
address $8t + 6$ deliver? And a word fetch from $8t + 5$? What happens with a
word store to $8t + 5$?
\end{exercise}

\begin{exercise}
\label{ex:step08-store}
A long \cc{0xDEADBEEF} is stored at byte address 12. Which tag and offset is
that, what are \cc{we00} to \cc{we11}, and which 16 bits does each SPRAM get at
\cc{DATAIN}?
\end{exercise}

\section{What the Shifters Cost}
\label{sec:step08-cost}

Two shifters, two extenders and a mask generator: how much hardware is that?

\outputfile{step08_report.txt}

The modules without interface ports appear in the list: the five \cc{quad_*}
modules together have $129 + 129 + 209 + 209 + 53 = 729$ cells, all of them
LUTs. That is more than the whole design of step 07, for logic that does
nothing but move bytes around. Each output bit of a shifter is a multiplexer
that chooses among up to eight input bits, controlled by three offset bits;
such a function has up to eleven inputs and needs several 4-input LUTs, and
there are 64 output bits.

These per-module numbers come from synthesising every module on its own, for
arbitrary inputs. In the real design, yosys flattens everything and removes
logic whose inputs are constant or whose outputs are unused. In this step the
loader and the debugger only access bytes, only the low byte of
\cc{data_in} is ever set, and only the low byte of \cc{data_out} is used. So the whole design has fewer LUTs than the five modules
alone:

\outputfile{step08_yosys_stat.txt}

To see what the four shift and extend units cost here, we replaced them by
plain wires (\cc{assign data_out = data_in;}) and synthesised again: 352 LUTs
instead of 696. So even in this step they account for half of the LUTs of the
design. In step 09, where the processor uses all access sizes and all 64 bits,
the same experiment gives 2843 against 2232 LUTs: there they cost about 600
LUTs, close to the $129 + 129 + 209 + 209 = 676$ of the report. The slowest path of the design, which limits it to 27.2\,MHz,
also runs through them: from the register \cc{op} of the loader, through the
switch, the selection of the data type and the shift and extend logic, to the
display.

The other number worth noticing: \cc{ICESTORM_SPRAM} 4 of 4. The memory of the
ULM takes all SPRAMs of the chip.

\section{Two Users, One Memory}
\label{sec:step08-switch}

\svfile{08_dev_loader/dev_ram_switch.sv}

The switch is a multiplexer for a whole interface. Its ports are modports
(Section~\ref{sec:step07-interface}), and their roles show the direction of the
wires: towards the loader and the debugger the switch is a \cc{server} (it
receives requests and returns \cc{data_out}), towards \cc{dev_ram} it is a
\cc{client}. The \cc{case} selects which user's requests go to the memory:
2-to-1 multiplexers for all bits of \cc{op},
\cc{data_type}, \cc{addr} and \cc{data_in}. The data from the memory goes to
the selected user, and the other one gets 0.

The debugger:

\svpart{firstline=12,lastline=51}{08_dev_loader/dev_ram_debugger.sv}

The two buttons go through the debouncer of step 07
(Section~\ref{sec:step07-devio}), so each press gives one pulse of one
cycle, however much the contacts bounce. \cc{BTN1} fetches address 0 and
sets the counter \cc{addr} to 1; each \cc{BTN2} fetches the address in
\cc{addr} and increments it. So the presses \cc{BTN1}, \cc{BTN2},
\cc{BTN2}, \dots{} read the addresses 0, 1, 2, \dots. The output \cc{out}
is the address, or with \cc{BTN3} the byte read, both extended to 64 bits.
\cc{LED4} lights if \cc{out} has bits set above the low byte, which the
display cannot show: without \cc{BTN3}, from address 256 on.

The directory also contains \cc{dev_io_switch.sv}, the same kind of switch for
\cc{if_io}. It is not used in this step; step 09 needs it.

\section{The Loader}
\label{sec:step08-loader}

The loader reads one character at a time from \cc{dev_io} through the
\cc{client} modport of \cc{if_io}. First, a combinational circuit decides
whether the character is a hexadecimal digit, and which value it has:

\svpart{firstline=18,lastline=49}{08_dev_loader/dev_loader.sv}

\cc{"0"} is a string literal of one character, and in a comparison it stands
for its 8-bit ASCII code, \cc{8'h30}. So the \cc{case} compares
\cc{io.getc_char} with 22 constants, and the block is a function of the 8 bits
of the character with 5 output bits: a small decoder, a handful of LUTs, not
a search through a table at run time. The first line gives the outputs their
values for all other characters, so there are no missing cases.

Then two nibbles become a byte:

\svpart{firstline=55,lastline=101}{08_dev_loader/dev_loader.sv}

The flip-flop \cc{nibble_sel} says which half comes next, 1 for the high
nibble. When a digit arrives, the loader pops it (\cc{io.getc_pop <= 1}),
stores its value in \cc{high_nibble} or \cc{low_nibble}, and toggles
\cc{nibble_sel}. After the low nibble, \cc{new_byte} is 1 for one cycle. A
character that is neither a digit nor \cc{0x04} is popped and ignored. The
character \cc{0x04} is ASCII ``end of transmission'', what a terminal sends for
Ctrl-D; it sets \cc{got_eof}, one cycle later \cc{done} is 1, and the loader
stops reading. The condition \cc{!io.getc_pop} keeps the loader from popping
in two consecutive cycles, as required by the FIFO
(Section~\ref{sec:step07-fifo}).

Finally, each new byte is stored at the next address:

\svpart{firstline=107,lastline=131}{08_dev_loader/dev_loader.sv}

\cc{ram.data_type} is set to \cc{RAM_BYTE} in the \cc{initial} block and never
changed, and only the low 8 bits of \cc{ram.data_in} are ever written: this is
the constant input that made the shifters cheaper in
Section~\ref{sec:step08-cost}.

How fast can the loader take characters? It decides at a clock edge to pop
(\cc{io.getc_pop <= 1}); in the next cycle the pop happens, and in the cycle
after that \cc{getc_en} is 0, because \cc{getc_char} still shows the popped
character (Section~\ref{sec:step07-devio}). So the loader takes at most one
character every three cycles, and never a stale one.

In the simulation (\cc{make sim}), the testbench delivers the characters
\cc{0} to \cc{9}, \cc{a}, \cc{b} and \cc{0x04}, one every two clock
cycles, in the same way as in step 07. That is faster than the loader takes
them, so several characters wait in the input buffer:

\cpppart{firstline=38,lastline=64}{08_dev_loader/tb_test.cpp}

In the waveform you see \cc{ram.op} at \cc{RAM_STORE} at time 28 with
address 0 and data \cc{0x01}, then every 12 time units (six clock cycles,
two characters) the next byte, up to address 5 with \cc{0xab} at time 88,
and \cc{done} becoming 1 at time 94. Then the testbench plays the user of
the debugger: it presses \cc{BTN1} once and \cc{BTN2} five times, each press
bouncing as in step 07, and after each press reads the address and, with
\cc{BTN3}, the byte from the debugger's output \cc{debugger_out}, which
\cc{test.sv} marks as \cc{/* verilator public */}:

\outputfile{step08_sim.txt}

Verilator warns about \cc{MULTIDRIVEN} signals in \cc{dev_loader.sv}: as in
step 02, these are the \cc{initial} blocks that set power-up values.

\begin{exercise}
\label{ex:step08-loader}
What does the loader store if you type \cc{3 2 1} (return) \cc{0} and then
Ctrl-D? What if you type \cc{123} and Ctrl-D? What does the display show, and
is \cc{LED4} on, before you press Ctrl-D in the second case?
\end{exercise}

\begin{pitfall}
A testbench that delivers characters at the speed of the real UART would
never have found the problem of the pitfall in Section~\ref{sec:step07-devio}:
at 9600 baud the loader has taken each character long before the next
arrives. With the earlier version of \cc{dev_io}, the testbench above shows
the loader storing \cc{00 12 34 56 78 9a} instead of \cc{01 23 45 67 89 ab}:
it reads the stale \cc{'0'} a second time, pops the \cc{'1'} without having
read it, and so on. The \cc{0x04} is popped unread as well, so the loader is
never done. Test at the limits, not only at the speed you expect.
\end{pitfall}

\section{A Declaration Is Not a Connection}
\label{sec:step08-rst}

The reset of the loader is wired in \cc{test.sv} in two lines:
\begin{lstlisting}
logic rst;
assign rst = BTN1 & BTN3;
\end{lstlisting}
An earlier version said it in one line,
\begin{lstlisting}
logic rst = BTN1 & BTN3;
\end{lstlisting}
which looks the same, but in SystemVerilog it is something else: \cc{=} in
the declaration of a variable gives it an \emph{initial value}, computed once
at the start, just like an \cc{initial} block (Chapter~\ref{ch:step02}).
After that, nothing drives \cc{rst}. A tiny test with Verilator:
\begin{lstlisting}
module tb;
  logic a = 0, b = 0;
  logic r = a & b;
  initial begin
    #1 a = 1; b = 1;
    #1 $display("a=%b b=%b r=%b", a, b, r);
    $finish;
  end
endmodule
\end{lstlisting}
\outputfile{step08_rinit.txt}
yosys interpreted the line differently. Its log for the earlier version said
\outputfile{step08_yosys_rst.txt}
which means it built combinational logic for \cc{rst}, as for an \cc{assign}.
So on the board the two buttons did reset the loader, but in the simulation
\cc{rst} stayed 0 forever, whatever the testbench did with the buttons.

\begin{pitfall}
Never write \cc{logic x = expression;} for a signal that should follow the
expression. Write \cc{logic x;} and \cc{assign x = expression;}. With a
constant, as in \cc{logic rst = 0;} of step 07, the declaration is fine: it is
a power-up value.
\end{pitfall}

\begin{exercise}
\label{ex:step08-rst}
Extend the testbench: after reading the memory, it presses \cc{BTN1} and
\cc{BTN3} together for five clock cycles, sets \cc{numLoaded} to 0 and
replaces the characters in \cc{prog} by \cc{fedcba987654}, so that they are
sent again, waits for \cc{LED3}, and reads the memory again with the
debugger. What do you expect? Why may the reset come straight from the
buttons, without a debouncer? And what would the simulation have shown with
the earlier declaration of \cc{rst}?
\end{exercise}

\section{What Is Next}

We now have everything around a processor: an I/O device with buffers, a
memory of 128\,KiB with byte, word, long and quad accesses, and a loader that
puts a program into it. What is missing is the processor itself: registers, an
ALU, and a control unit that fetches instructions from the memory, decodes
and executes them. That is step 09, the ULM.

\solutions

\solution
Offset 6: shifted left by 48 bits, the quad is \cc{0xCDEF000000000000}, and
the top 16 bits zero-extended give \cc{0x000000000000CDEF}. Offset 5: shifted
by 40 bits, \cc{0xABCDEF0000000000}, so the result is \cc{0xABCD}: unaligned
fetches deliver the two bytes at addresses $8t+5$ and $8t+6$, as long as they
are in the same quad. An unaligned word store writes nothing, since
\cc{quad_we_mask} has no case for a word at offset 5 and leaves the mask 0.

\solution
$12 = 8 \cdot 1 + 4$: tag 1, offset 4. The mask is \cc{16'h00FF}, so
\cc{we00} = \cc{we01} = \cc{4'b0000} and \cc{we10} = \cc{we11} =
\cc{4'b1111}. \cc{quad_ext_low} gives \cc{0xDEADBEEF00000000}, \cc{quad_rshift}
by 32 bits gives \cc{0x00000000DEADBEEF}. So \cc{spram00} and \cc{spram01} get
\cc{0x0000} (and do not write it), \cc{spram10} gets \cc{0xDEAD} and
\cc{spram11} gets \cc{0xBEEF}, both written at address 1.

\solution
Spaces and the return character are skipped, so \cc{3 2 1}, return, \cc{0}
stores \cc{0x32} at address 0 and \cc{0x10} at address 1, just like
\cc{3210}. With \cc{123}, the loader stores \cc{0x12} at address 0. The
\cc{3} becomes the high nibble of the next byte: \cc{low_nibble} is set to 0,
so the display shows \cc{30}, and since \cc{nibble_sel} is 0, \cc{byte_valid}
is 0 and \cc{LED4} is on. Ctrl-D ends the loading anyway; the incomplete byte
is never stored.

\solution
\cc{LED3} goes off with the reset, the loader is done again 43 cycles after
the reset, and the debugger reads \cc{fe dc ba 98 76 54} at the
addresses 0 to 5: the reset cleared \cc{done}, \cc{nibble_sel} and the
loader's \cc{addr}, so the stores started at address 0 again. The five
cycles of \cc{BTN1} are far too short for the debouncer of the debugger, so
it does not fetch. A reset needs no debouncer because it is idempotent
(Chapter~\ref{ch:step04}): resetting several times in a row, as a bouncing
button would, has the same effect as resetting once. With the earlier
declaration, \cc{rst} would have stayed 0 in the simulation: \cc{LED3} would
stay on, no characters would be taken, and the debugger would read the old
bytes \cc{01 23 45 67 89 ab}.
